\scanpage{94}

$Q'$: clôture tot. réelle de $Q$ \ldots

\margin{$\alpha=2$}

La condition a équivaut à la suivante, si $q\in Q^{*\unsure{+}}$ est $>0$:

\diagram{\begin{tikzcd}[row sep=large,cells={nodes={anchor=west}}]
\parbox{.88\linewidth}{$\left\{\begin{array}{l}\deleted{\unsure{M\subset M}}\ \inserted{\text{Pour tout }\lambda\in M,\ \exists r>0\text{ tel que}}\\\unsure{\lambda}^r\in Q'(i)^*,\quad\text{et pour tout }\lambda\in M,\text{ il existe un}\\\text{entier }i\text{ tel que }|\lambda|=q^{i/2}\quad\text{i.e. }\lambda\bar\lambda=q^i\end{array}\right.$\par\margin{N.B. Le $i$ en question est uniquement déterminé et n'est autre que $j^*(\lambda)$.}} \arrow[d,Leftrightarrow] \\
\parbox{.88\linewidth}{$\left\{\begin{array}{l}\text{Pour tout }\lambda\in M,\ \deleted{\text{il existe}}\ \text{on a }|\lambda|=q^{j^*(\lambda)/2}\quad\text{i.e. }\lambda\bar\lambda=q^{j^*(\lambda)}\end{array}\right.$}
\end{tikzcd}}

\uline{Cor} Soit $\lambda\in\overline Q$, tel que \deleted{il existe $q\in Q^{*+}$, avec}
\[
\lambda_\alpha\bar\lambda_\alpha=\lambda\bar\lambda\in Q
\]
\inserted{$\lambda\bar\lambda=q$}
\inserted{$|\lambda_\alpha|^2=|\lambda|^2$}
pour tout conjugué $\lambda_\alpha$ de $\lambda$. Alors $\exists r>0$, tel que $\lambda^r\in Q'(i)$, [donc la condition ne dépend pas des plongements de $\unsure{Q(\lambda)|R(\lambda)}$ dans $\overline Q\subset\mathbb C$], et il existe $M,\varepsilon,j\ldots$ \inserted{relatif à $q$} tel que $\lambda\in M$ (et réciproquement).

\uline{Proposition} Soit $\lambda\in\overline Q^*$, et $q\in Q^{++}$. Conditions équivalentes

(i) $\exists$ entier $i$ tel que pour $\forall$ conjugué $\lambda_\alpha$ de $\lambda$ on ait $|\lambda_\alpha|=q^{i/2}$.
\inserted{i.e. $\lambda_\alpha\bar\lambda_\alpha=q^i$}

\inserted{(\unsure{ii bis}) $\exists$ entier $n>0$ tel que $\lambda^n\in Q'(i)$, et \deleted{\unsure{$\exists$}} entier $i$ tel que $\lambda^n\overline{\lambda^n}=q^{in}$.}

(ii) $\lambda\in Q'(i)$, il existe entier $i$ tel que $|\lambda|=q^{i/2}$ i.e. $\lambda\bar\lambda=q^i$.
\inserted{$|\lambda|=q^{i/2}$, i.e. $\lambda\bar\lambda=q^i$}

(iii) $\exists$ un groupe de t.m. $G$ sur $Q$, et un élément $f$ de $G(Q)$,
\diagram{\begin{tikzcd}G_m\arrow[r,"j"]&G\arrow[r,"\varepsilon"]&G_m\end{tikzcd}},
avec $\varepsilon j(\lambda)=\lambda^2$, \deleted{tel que $\unsure{\operatorname{Ker}\varepsilon}$} \inserted{avec $\operatorname{Ker}\varepsilon$ \uline{compact}}, $\varepsilon(f)=q$, tels que $\lambda$ soit une valeur propre de $f$ pour un $G$-module de dim finie \inserted{sur $Q$} convenable, i.e. $\lambda$ appartient à l'image de l'hom.
\[
M=D(G_{\overline Q})\longrightarrow\overline Q^*
\]
défini par $f$.

(iv) Il existe une algèbre de \unsure{Riemann} \unsure{commutative} sur $Q$ et $f\in A^*$ \inserted{et $i\unsure{\in}\mathbb N$, tels que $ff'=q^i$} et un $Q$-hom. $\unsure{A}\longrightarrow\overline Q$, appliquant $f$ dans $\lambda$.

\margin{93}
\margin{\uline{N.B.} Si un tel $\lambda$ est réel (donc tot réel!) on a $\lambda^2=q^i$ donc $\lambda=(q^{1/2})^i$, \uline{cas trivial}}

\margin{\deleted{N.B. La condition \uline{n'implique pas} que $\lambda\in Q'(i)^*$, comme on voit p. ex. en prenant $\lambda=p^{1/4}$; Alors $Q(\lambda)$ n'est ni tot. imaginaire ($Q(p^{1/4})$ est réel) ni tot réel (car il est conjugué à $Q(\sqrt{-\sqrt p})$ qui est \unsure{purement imaginaire}, donc n'est pas $\subset Q'(i)$)}}
